\section{练习}

\begin{question}\begin{drill}[半厚度填充与半面积填充比较]\label{e5:drill:1}
同一真空电容器 $C_0$，分别用介质 $\kappa$ 填充一半面积并贯穿间隙，或填充全部面积的一半厚度。求电容并比较。
\end{drill}
\typeout{ECFIG:fig:e5:drill:1:s}\fig{fig:e5:drill:1:s}{题设：半厚度填充与半面积填充比较}{\begin{tikzpicture}[scale=.88]
\draw[world](0,1.25)--(4.5,1.25);\node[right]at(4.5,1.25){$A$};\draw[world](0,0.0)--(4.5,0.0);\node[right]at(4.5,0.0){$B$};\fill[blue!12](0,0)rectangle(2.2,1.25);\draw[guide](2.2,0)--(2.2,1.25);\node at(1,.6){$\kappa$};\node at(3.4,.6){真空};\begin{scope}[xshift=7cm]\draw[world](0,1.25)--(4.5,1.25);\node[right]at(4.5,1.25){$A$};\draw[world](0,0.0)--(4.5,0.0);\node[right]at(4.5,0.0){$B$};\fill[blue!12](0,0)rectangle(4.5,.65);\draw[guide](0,.65)--(4.5,.65);\node at(2,.3){$\kappa$};\node at(2,.95){真空};\end{scope}
\end{tikzpicture}}
\end{question}
\answerlink{e5:drill:1}
\begin{question}\begin{drill}[两只孤立并联电容的电荷转移]\label{e5:drill:2}
两只电容原为 $C_1,C_2$，并联充至 $U_0$ 后整体断开电源但保持并联。把第一只填满介质使之为 $\kappa C_1$。求末电压及两只电量变化。
\end{drill}
\typeout{ECFIG:fig:e5:drill:2:s}\fig{fig:e5:drill:2:s}{两只原电压同为 $U_0$，整个并联组与电源断开}{\begin{tikzpicture}
\draw(0,0)to[C,l_={$C_1$}](0,3)--(5,3)to[C,l={$C_2$}](5,0)--(0,0);\node at(2.5,3.7){原上端总量 $(C_1+C_2)U_0$};\node at(2.5,-.7){插入时保留两条连线};
\end{tikzpicture}}
\end{question}
\answerlink{e5:drill:2}
\begin{question}\begin{drill}[介质中的电击穿]\label{e5:drill:3}
两层介质厚度 $d_1,d_2$、介电常数 $\varepsilon_1,\varepsilon_2$，击穿场分别 $E_{b1},E_{b2}$。无界面自由电荷，求最大安全电压。
\end{drill}
\typeout{ECFIG:fig:e5:drill:3:s}\fig{fig:e5:drill:3:s}{两层厚度与各自击穿参数}{\begin{tikzpicture}

\fill[blue!10](0,1.3)rectangle(5,3);\fill[orange!10](0,0)rectangle(5,1.3);\draw[world](0,0)--(5,0)(0,3)--(5,3);\draw[guide](0,1.3)--(5,1.3);\node at(1.4,2.2){$\varepsilon_1,E_{b1}$};\node at(1.4,.6){$\varepsilon_2,E_{b2}$};\draw[dim](-.6,1.3)--(-.6,3)node[midway,left]{$d_1$};\draw[dim](-.6,0)--(-.6,1.3)node[midway,left]{$d_2$};

\end{tikzpicture}}
\end{question}
\answerlink{e5:drill:3}
\begin{question}\begin{drill}[介质球壳填充的球形电容]\label{e5:drill:4}
内外金属球半径 $a,b$，介质 $\kappa$ 填在 $a<r<c$，其余为真空。求电容及 $r=c$ 介质端面的束缚密度，内球电量 $+Q$。
\end{drill}
\typeout{ECFIG:fig:e5:drill:4:s}\fig{fig:e5:drill:4:s}{介质只占内球外的 $a<r<c$，外段仍是真空}{\begin{tikzpicture}

\fill[blue!13](0,0)circle(1.3);\fill[gray!15](0,0)circle(.65);\draw[world](0,0)circle(.65);\draw[guide](0,0)circle(1.3);\draw[world](0,0)circle(2);\draw[dim](0,0)--(.46,.46)node[midway,above left]{$a$};\draw[dim](0,0)--(-.92,.92)node[pos=.8,above right]{$c$};\draw[dim](0,0)--(1.41,-1.41)node[pos=.85,above right]{$b$};\node at(-.95,-.2){$\kappa$};\draw(-1.6,-.6)--(-2.7,-1.2)node[left]{真空};
\node[right]at(2.7,.4){内球 $+Q$};
\end{tikzpicture}}
\end{question}
\answerlink{e5:drill:4}
\begin{question}\begin{drill}[导体片横向部分插入]\label{e5:drill:5}
矩形极板长 $L$、宽 $b$、间距 $d$，厚 $t<d$ 的中性金属片横向插入长度 $x$，不接触两极。求电容和恒压 $U$ 下的吸入力。
\end{drill}
\typeout{ECFIG:fig:e5:drill:5:s}\fig{fig:e5:drill:5:s}{金属片厚度只占部分间隙，仍与两极绝缘}{\begin{tikzpicture}

\draw[world](0,0)--(5,0)(0,2.4)--(5,2.4);\fill[gray!25](-1,.8)rectangle(2,1.6);\draw[world](-1,.8)rectangle(2,1.6);\draw[dim](0,-.5)--(2,-.5)node[midway,below]{$x$};\draw[dim](0,3)--(5,3)node[midway,above]{$L$};\draw[dim](5.6,0)--(5.6,2.4)node[midway,right]{$d$};\draw[dim](-1.6,.8)--(-1.6,1.6)node[midway,left]{$t$};\node at(.5,1.2){金属};\node at(3.7,1.2){真空};

\end{tikzpicture}}
\end{question}
\answerlink{e5:drill:5}
\begin{question}\begin{drill}[介质吸入后获得的速度]\label{e5:drill:6}
介质块质量 $M$、相对介电常数 $\kappa$，从刚进入平行板处静止释放，完全填入后电容由 $C_0$ 变为 $\kappa C_0$。忽略摩擦与损耗。分别求断电定量 $Q$ 和接恒压 $U$ 时的末速度。
\end{drill}
\typeout{ECFIG:fig:e5:drill:6:s}\fig{fig:e5:drill:6:s}{题设：介质吸入后获得的速度}{\begin{tikzpicture}[scale=.88]
\draw[world](0,1.25)--(4.5,1.25);\node[right]at(4.5,1.25){$A$};\draw[world](0,0.0)--(4.5,0.0);\node[right]at(4.5,0.0){$B$};\fill[blue!12](-2,0)rectangle(0,1.25);\draw[world](-2,0)rectangle(0,1.25);\node at(-1,.6){$M,\kappa$};\draw[flow](-1,1.8)--(.5,1.8);\node at(2.5,.6){初始真空间隙};
\end{tikzpicture}}
\end{question}
\answerlink{e5:drill:6}
\begin{question}\begin{drill}[界面电场方向的折转]\label{e5:drill:7}
两均匀介质界面无自由电荷。场在介质1中与法线夹角 $\theta_1$，在介质2中夹角 $\theta_2$，介电常数为 $\varepsilon_1,\varepsilon_2$。求角度关系并判断进入高介电常数侧时的折转。
\end{drill}
\typeout{ECFIG:fig:e5:drill:7:s}\fig{fig:e5:drill:7:s}{两个角都相对法线，图示进入较高介电常数一侧}{\begin{tikzpicture}

\fill[blue!10](-3,-2)rectangle(3,0);\draw[world](-3,0)--(3,0);\draw[guide](0,-2.2)--(0,2.3);\draw[flow](-1.2,2)--(-.05,.083);\draw[flow](0,0)--(2,-1.4);\draw(0,.7)arc(90:121:.7);\node at(-.35,.95){$\theta_1$};\draw(0,-.7)arc(-90:-35:.7);\node at(.55,-.92){$\theta_2$};\node at(-2,1){$\varepsilon_1$};\node at(-2,-1){$\varepsilon_2$};\node[right]at(.1,2){法线};

\end{tikzpicture}}
\end{question}
\answerlink{e5:drill:7}
\begin{question}\begin{drill}[介质板拉出所需外功]\label{e5:drill:8}
真空电容为 $C_0$ 的电容器已完全填入介质 $\kappa$，初电压 $U_0$。分别在此时断电和始终接恒压源两种情况下，缓慢把介质完全拉出，求外功。
\end{drill}
\typeout{ECFIG:fig:e5:drill:8:s}\fig{fig:e5:drill:8:s}{题设：介质板拉出所需外功}{\begin{tikzpicture}[scale=.88]
\fill[blue!12](0,0)rectangle(4.5,1.25);\draw[world](0,1.25)--(4.5,1.25);\node[right]at(4.5,1.25){$A$};\draw[world](0,0.0)--(4.5,0.0);\node[right]at(4.5,0.0){$B$};\node at(2.25,.6){$\kappa$};\draw[force](0,.6)--(-1.5,.6)node[left]{拉出};
\end{tikzpicture}}
\end{question}
\answerlink{e5:drill:8}
\begin{question}\begin{drill}[点电荷外的同心介质壳]\label{e5:drill:9}
点电荷 $q$ 位于球心，$a<r<b$ 有介电常数 $\kappa\varepsilon_0$ 的球壳，其余真空。求三段电场、壳两面的束缚电量及介质在球心产生的反作用电势。
\end{drill}
\typeout{ECFIG:fig:e5:drill:9:s}\fig{fig:e5:drill:9:s}{球心点电荷与两侧真空中的同心介质壳}{\begin{tikzpicture}

\fill[blue!12](0,0)circle(2);\fill[white](0,0)circle(1.1);\draw[world](0,0)circle(1.1);\draw[world](0,0)circle(2);\fill(0,0)circle(2pt)node[below left]{$q$};\draw[dim](0,0)--(-.78,.78)node[midway,above right]{$a$};\draw[dim](0,0)--(1.41,-1.41)node[pos=.8,above right]{$b$};\node at(-1.55,-.4){$\kappa$};

\end{tikzpicture}}
\end{question}
\answerlink{e5:drill:9}
